% Options for packages loaded elsewhere
\PassOptionsToPackage{unicode}{hyperref}
\PassOptionsToPackage{hyphens}{url}
\PassOptionsToPackage{dvipsnames,svgnames,x11names}{xcolor}
\documentclass[
  11pt,
]{article}
\usepackage{xcolor}
\usepackage[margin=27mm]{geometry}
\usepackage{amsmath,amssymb}
\setcounter{secnumdepth}{-\maxdimen} % remove section numbering
\usepackage{iftex}
\ifPDFTeX
  \usepackage[T1]{fontenc}
  \usepackage[utf8]{inputenc}
  \usepackage{textcomp} % provide euro and other symbols
\else % if luatex or xetex
  \usepackage{unicode-math} % this also loads fontspec
  \defaultfontfeatures{Scale=MatchLowercase}
  \defaultfontfeatures[\rmfamily]{Ligatures=TeX,Scale=1}
\fi
\usepackage{lmodern}
\ifPDFTeX\else
  % xetex/luatex font selection
  \setmainfont[]{Cambria}
  \setmonofont[]{Consolas}
\fi
% Use upquote if available, for straight quotes in verbatim environments
\IfFileExists{upquote.sty}{\usepackage{upquote}}{}
\IfFileExists{microtype.sty}{% use microtype if available
  \usepackage[]{microtype}
  \UseMicrotypeSet[protrusion]{basicmath} % disable protrusion for tt fonts
}{}
\makeatletter
\@ifundefined{KOMAClassName}{% if non-KOMA class
  \IfFileExists{parskip.sty}{%
    \usepackage{parskip}
  }{% else
    \setlength{\parindent}{0pt}
    \setlength{\parskip}{6pt plus 2pt minus 1pt}}
}{% if KOMA class
  \KOMAoptions{parskip=half}}
\makeatother
\usepackage{longtable,booktabs,array}
\newcounter{none} % for unnumbered tables
\usepackage{calc} % for calculating minipage widths
% Correct order of tables after \paragraph or \subparagraph
\usepackage{etoolbox}
\makeatletter
\patchcmd\longtable{\par}{\if@noskipsec\mbox{}\fi\par}{}{}
\makeatother
% Allow footnotes in longtable head/foot
\IfFileExists{footnotehyper.sty}{\usepackage{footnotehyper}}{\usepackage{footnote}}
\makesavenoteenv{longtable}
\setlength{\emergencystretch}{3em} % prevent overfull lines
\providecommand{\tightlist}{%
  \setlength{\itemsep}{0pt}\setlength{\parskip}{0pt}}
\usepackage{bookmark}
\IfFileExists{xurl.sty}{\usepackage{xurl}}{} % add URL line breaks if available
\urlstyle{same}
\hypersetup{
  pdftitle={WP-14 Literature Review --- Adversarial Novelty Adjudication (Run 02)},
  pdfauthor={The Privacy-is-Value Research Programme},
  colorlinks=true,
  linkcolor={black},
  filecolor={Maroon},
  citecolor={Blue},
  urlcolor={blue},
  pdfcreator={LaTeX via pandoc}}

\title{WP-14 Literature Review --- Adversarial Novelty Adjudication (Run
02)}
\author{The Privacy-is-Value Research Programme}
\date{2026-07-15}

\begin{document}
\maketitle

\section{Lit-review runtime · run
manifest}\label{lit-review-runtime-run-manifest}

\textbf{Runtime:} \texttt{weis-litreview-runtime} (V6 pipeline workflow)
\textbf{Brief:}
\texttt{chronicles/BRIEF\_litreview\_runtime\_WEIS\_2026-07-14.md}
(Matsuo venue+gap calls; C17; defects D2/D3/D4b in scope)
\textbf{Purpose:} produce a \emph{defensible} novelty claim for WP-14
(WEIS 2027), not a summary. Mitchell runs the loop and owns all
conjecture assignment; the runtime produces candidates and structure
only (brief §2, §9). \textbf{Run ID:} wf\_ba400906-67b \textbf{Date:}
2026-07-15 \textbf{Operator role:} A0 (orchestrator), extending A10
(prior-art scout) with an adversarial dual-seat loop.

\begin{center}\rule{0.5\linewidth}{0.5pt}\end{center}

\subsection{Seat configuration (D4b auditability, brief
§5)}\label{seat-configuration-d4b-auditability-brief-5}

Seat separation is \textbf{structural, not prompt-only} (D3). Each seat
is a distinct workflow \texttt{agent()} invocation with its own isolated
context; no seat reads another seat's reasoning trace.

{\def\LTcaptype{none} % do not increment counter
\begin{longtable}[]{@{}
  >{\raggedright\arraybackslash}p{(\linewidth - 6\tabcolsep) * \real{0.2500}}
  >{\raggedright\arraybackslash}p{(\linewidth - 6\tabcolsep) * \real{0.2500}}
  >{\raggedright\arraybackslash}p{(\linewidth - 6\tabcolsep) * \real{0.2500}}
  >{\raggedright\arraybackslash}p{(\linewidth - 6\tabcolsep) * \real{0.2500}}@{}}
\toprule\noalign{}
\begin{minipage}[b]{\linewidth}\raggedright
Seat
\end{minipage} & \begin{minipage}[b]{\linewidth}\raggedright
Role
\end{minipage} & \begin{minipage}[b]{\linewidth}\raggedright
Isolation
\end{minipage} & \begin{minipage}[b]{\linewidth}\raggedright
Sees
\end{minipage} \\
\midrule\noalign{}
\endhead
\bottomrule\noalign{}
\endlastfoot
extract:* & per-paper schema extraction (one per corpus item) & isolated
invocation per paper & only its own paper hint + the PVM position +
WebSearch \\
prove:* & PROVER --- builds the case that the candidate is novel &
isolated invocation & the candidate claim + corpus digest + WebSearch \\
refute:* & DELEGATOR --- hunts covering prior art to destroy the claim &
isolated invocation, \textbf{does not see the prover's argument} & the
candidate claim + corpus digest + WebSearch \\
judge:* & JUDGE --- rules VALIDATED / MIRAGE / BLOCKED & isolated
invocation & the candidate claim + both seats' structured outputs \\
synthesise:* & assembles the four outputs & isolated invocation & the
full corpus + all verdicts \\
\end{longtable}
}

\textbf{D4b note (honest limit).} All seats run on the same underlying
model. The separation this runtime enforces is \emph{context isolation}
(each seat forms its own search and reasoning with no shared trace), not
distinct model weights. The delegator is given the full corpus and
independent WebSearch capacity and is instructed to default to finding
prior art; the adversariality is real but is seat-context-level, not
weights-level. A run that returns all-VALIDATED is treated as a
\textbf{failed} enforcement of D2/D3, not a success --- MIRAGE is the
expected, valuable verdict.

\subsection{Hardening constraints
enforced}\label{hardening-constraints-enforced}

\begin{itemize}
\tightlist
\item
  \textbf{D2 (no absence-as-novelty):} the delegator must report absence
  as \texttt{not\_found\_in\_corpus(n,\ range)} with web scope, never as
  \texttt{novel}; every contribution claim carries a verdict + the
  corpus tested + a confidence label.
\item
  \textbf{D3 (structural isolation):} prover and delegator are separate
  invocations; the delegator never sees the prover's case.
\item
  \textbf{D4b (seat separation auditable):} this manifest records the
  seat config; see the honest limit above.
\item
  \textbf{C81 hold-back:} the Privacy Pools V2 constraint-reduction
  result does not enter the runtime or any output.
\end{itemize}

\subsection{Corpus (Tier A, challenge
corpus)}\label{corpus-tier-a-challenge-corpus}

20 items across five strands: propertization (6), data-as-labour (2),
WEIS-native economics-of-privacy (7), privacy quantification (3),
standards/PIA (2). Provenance-gated: each carries a DOI/URL and a
\texttt{provenance\_confidence}; \texttt{uncertain} items are flagged
for first-person verification in \texttt{bibliography.md}.

\subsection{Outputs (this directory)}\label{outputs-this-directory}

\begin{itemize}
\tightlist
\item
  \texttt{gap\_table.md} --- corpus clustered on unit\_of\_analysis ×
  economic\_framing; the \texttt{system\_property\ ×\ market\_pricing}
  cell called out.
\item
  \texttt{contribution\_claims.md} --- CTR-shaped novelty candidates,
  each with verdict, residual novel core, corpus tested, confidence.
\item
  \texttt{ctr\_candidates.md} --- conjecture-shaped residues (CTR-series
  only; assignment is the First Person's).
\item
  \texttt{bibliography.md} --- full citations with provenance,
  WEIS-submission ready.
\end{itemize}

\emph{Generated artifacts; do not hand-edit. Fix the runtime inputs and
rebuild.}

\newpage

\section{Gap table: the privacy-value corpus on unit-of-analysis x
economic-framing}\label{gap-table-the-privacy-value-corpus-on-unit-of-analysis-x-economic-framing}

\subsection{Gap grid: unit\_of\_analysis x
economic\_framing}\label{gap-grid-unit_of_analysis-x-economic_framing}

Corpus n = 20. Each cell lists the corpus items that land there. Short
keys resolve to full citations in the bibliography.

{\def\LTcaptype{none} % do not increment counter
\begin{longtable}[]{@{}
  >{\raggedright\arraybackslash}p{(\linewidth - 8\tabcolsep) * \real{0.2000}}
  >{\raggedright\arraybackslash}p{(\linewidth - 8\tabcolsep) * \real{0.2000}}
  >{\raggedright\arraybackslash}p{(\linewidth - 8\tabcolsep) * \real{0.2000}}
  >{\raggedright\arraybackslash}p{(\linewidth - 8\tabcolsep) * \real{0.2000}}
  >{\raggedright\arraybackslash}p{(\linewidth - 8\tabcolsep) * \real{0.2000}}@{}}
\toprule\noalign{}
\begin{minipage}[b]{\linewidth}\raggedright
unit ~framing
\end{minipage} & \begin{minipage}[b]{\linewidth}\raggedright
none
\end{minipage} & \begin{minipage}[b]{\linewidth}\raggedright
cost\_model
\end{minipage} & \begin{minipage}[b]{\linewidth}\raggedright
incentive\_model
\end{minipage} & \begin{minipage}[b]{\linewidth}\raggedright
market\_pricing
\end{minipage} \\
\midrule\noalign{}
\endhead
\bottomrule\noalign{}
\endlastfoot
\textbf{subject\_impact} & ISO/IEC 29134:2017; ISO/IEC 29100:2011 &
(empty) & Delacroix-Lawrence 2019 & Laudon 1996; Samuelson 2000;
Schwartz 2004; Purtova 2015/17; Prins 2006; Arrieta-Ibarra et al.~2018;
Posner-Weyl 2018; Acquisti-John-Loewenstein 2013; Collis et al.~2022;
Beresford-Kuebler-Preibusch 2012; Grossklags-Acquisti 2007 \\
\textbf{system\_property} & Dwork 2006; Sweeney 2002; Smith 2009 &
(empty) & Anderson 2001 & \textbf{(empty, 0 items)} \\
\textbf{both} & (empty) & (empty) & Acquisti-Taylor-Wagman 2016 &
Spiekermann et al.~2015 \\
\textbf{unclear} & (empty) & (empty) & (empty) & (empty) \\
\end{longtable}
}

Row/column totals: subject\_impact = 14, system\_property = 4, both = 2,
unclear = 0. Column totals: none = 5, cost\_model = 0, incentive\_model
= 3, market\_pricing = 12.

\subsection{The load-bearing cell: \{system\_property x
market\_pricing\}}\label{the-load-bearing-cell-system_property-x-market_pricing}

\textbf{Exact membership count in the corpus: 0.}

No item in this 20-paper corpus both (a) takes the unit of measurement
to be a property of the release mechanism or system, rather than the
reported impact on a named subject, and (b) prices that property as a
market quantity. The four system\_property items either carry no
economic framing at all (Dwork, Sweeney, Smith are formal leakage or
indistinguishability bounds with \texttt{economic\_framing\ =\ none}) or
frame economics as incentive misalignment rather than pricing (Anderson,
\texttt{incentive\_model}). Conversely, every market\_pricing item
measures a subject-reported valuation, a legal entitlement of a named
subject, or a market-position quantity attached to a subject, not a
system property.

\textbf{What this means for the WEIS contribution claim.} This emptiness
is a genuine structural feature of the assembled corpus, and it
motivates the PVM position, but under discipline D2 it must be reported
as
\texttt{not\_found\_in\_corpus(system\_property\ x\ market\_pricing,\ this\ 20-item\ corpus)}
and NOT as ``no prior work prices a system property.'' The adjudicated
verdict on CTR-LR-02 is explicit that the literal in-corpus emptiness is
a corpus-selection artefact: the external differential-privacy-markets
strand (Ghosh-Roth 2011 and successors) does price the DP disclosure
bound as a traded commodity and therefore occupies this cell outside the
corpus window. Confidence that the cell is empty \emph{within this
corpus}: high (certain by inspection). Confidence that the cell is empty
\emph{in the field}: refuted (see CTR-LR-02 verdict). The defensible
WEIS framing is therefore narrower than ``we are first to price a system
property'': it is ``we price a finite, adversary-relative, time-drifting
reconstruction bound held as inalienable subject-owned capital,'' which
is the residual that survives adjudication.

\subsection{Second grid: adversary\_model
coverage}\label{second-grid-adversary_model-coverage}

{\def\LTcaptype{none} % do not increment counter
\begin{longtable}[]{@{}
  >{\raggedright\arraybackslash}p{(\linewidth - 4\tabcolsep) * \real{0.3333}}
  >{\raggedright\arraybackslash}p{(\linewidth - 4\tabcolsep) * \real{0.3333}}
  >{\raggedright\arraybackslash}p{(\linewidth - 4\tabcolsep) * \real{0.3333}}@{}}
\toprule\noalign{}
\begin{minipage}[b]{\linewidth}\raggedright
adversary\_model
\end{minipage} & \begin{minipage}[b]{\linewidth}\raggedright
count
\end{minipage} & \begin{minipage}[b]{\linewidth}\raggedright
items
\end{minipage} \\
\midrule\noalign{}
\endhead
\bottomrule\noalign{}
\endlastfoot
\textbf{explicit} & 2 & Dwork 2006; Smith 2009 \\
\textbf{implicit} & 7 & Samuelson 2000; Delacroix-Lawrence 2019;
Acquisti-Taylor-Wagman 2016; Spiekermann et al.~2015; Anderson 2001;
Sweeney 2002; ISO/IEC 29134:2017 \\
\textbf{absent} & 11 & Laudon 1996; Schwartz 2004; Purtova 2015/17;
Prins 2006; Arrieta-Ibarra et al.~2018; Posner-Weyl 2018;
Acquisti-John-Loewenstein 2013; Collis et al.~2022;
Beresford-Kuebler-Preibusch 2012; Grossklags-Acquisti 2007; ISO/IEC
29100:2011 \\
\end{longtable}
}

Reading: an explicit adversary appears only in the two pure
formal-metric papers (Dwork, Smith), both of which have
\texttt{economic\_framing\ =\ none}. Every market\_pricing item has an
absent or (in two survey cases) merely implicit adversary. The
intersection \{explicit adversary\} x \{market\_pricing\} is also empty
in the corpus. This reinforces the same gap from a different axis:
adversary-relativity and pricing do not co-occur in any corpus item,
which is exactly the join the PVM reconstruction bound R(t) attempts,
and exactly the join that the CTR-LR-05 residual (adversary-relative
drifting horizon carried into a capital valuation) is scoped to.

\newpage

\section{Contribution claims: adversarially adjudicated novelty
candidates}\label{contribution-claims-adversarially-adjudicated-novelty-candidates}

\subsection{Novelty candidates: verdicts and surviving
residues}\label{novelty-candidates-verdicts-and-surviving-residues}

Five novelty candidates were tested against the assembled corpus plus
the adjudicators' web scope. \textbf{All five returned MIRAGE at the
altitude stated.} None returned VALIDATED and none returned BLOCKED.
This is the honest headline: the strongest form of each claim as written
is anticipated. What survives in each case is a narrower residual novel
core. The residues are ranked below by the adjudicated confidence that
the residual itself survives, strongest first.

The verdict vocabulary: VALIDATED = the stated claim survives as novel;
MIRAGE = the stated claim is anticipated but a narrower residue
survives; BLOCKED = withheld (for example under an existence-leak hold).
No claim here is VALIDATED, so every entry below is a
MIRAGE-with-residue.

\subsubsection{1. CTR-LR-02 (strongest surviving residue) â€''
MIRAGE}\label{ctr-lr-02-strongest-surviving-residue-uxe2-mirage}

\textbf{Claim as stated.} No prior work prices privacy as a market
quantity attached to a system property (a reconstruction bound) rather
than to subject-reported impact or to a technical metric with no market;
the \{system\_property x market\_pricing\} cell is empty.

\textbf{Why it MIRAGE'd.} The differential-privacy-markets strand
(Ghosh-Roth 2011 \emph{Selling Privacy at Auction}, Li-Miklau-Roth-Suciu
2013/14, Fleischer-Lyu 2012, Cummings et al.~2015) already prices the DP
disclosure bound epsilon, an adversary-independent system property, as a
traded commodity with per-subject compensation and a formal equilibrium
price. That occupies the asserted-empty cell and refutes the ``no
market'' framing.

\textbf{Residual novel core (survives).} PVM is the first, against the
surfaced corpus plus the DP-markets strand, to price a \emph{finite,
adversary-relative, time-drifting} reconstruction bound R(t), not a
static adversary-independent DP or accuracy budget, held as inalienable
subject-owned durable capital, with the observer-to-subject value gap
decomposed as market-position rent rather than super-additive
aggregation. It is the \emph{joint} quadruple (drift +
adversary-relativity + structural inalienability + rent decomposition)
that is uncovered.

\textbf{Corpus tested against.} 20-item corpus (system\_property tags:
Anderson, Dwork, Sweeney, Smith; market\_pricing tags: Laudon,
Posner-Weyl, Acquisti et al.) plus external DP-markets strand.

\textbf{Confidence that the residual survives: high (86\%).}

\subsubsection{2. CTR-LR-03 â€'' MIRAGE}\label{ctr-lr-03-uxe2-mirage}

\textbf{Claim as stated.} The subject-side near-zero surplus share is
derived as the subgame-perfect equilibrium of a consent-interface
bargaining game (notice-and-consent read as an ultimatum extensive
form), and IEEE 7012 propose-and-respond inverts the proposer role to
move that share off the floor.

\textbf{Why it MIRAGE'd.} Half A (near-zero responder share as the SPE
of a one-shot ultimatum) is the standard Gueth 1982 / Stahl-Rubinstein
result. Half B (IEEE 7012 / MyTerms inverting the proposer to shift
bargaining power) is the MyTerms standard's own self-description. The
bridge (reading the consent screen as the extensive form) is already
published: privacy-law scholarship (Norton 2016; Richards-Hartzog 2025)
maps notice-and-consent adhesion contracts onto take-it-or-leave-it
structures with the firm as unilateral proposer. The composite was not
found as one unit, but each half is separately covered and the join is
routine.

\textbf{Residual novel core (survives).} Casting the empirically
measured WTA/WTP floor and the MyTerms remedy as two evaluations of one
preference-free SPE share function of proposer identity, with the
subject-side share written as a structural venue-exclusion parameter s =
h\_X (the floor arising because the subject is excluded from the
clearing venue, a structural rather than cognitive origin), and IEEE
7012 as the specific proposer permutation carrying h\_X to an interior
share.

\textbf{Corpus tested against.} 20-item corpus plus ultimatum/Rubinstein
microeconomics, the WTA/WTP strand, the propertisation strand,
notice-and-consent-as-adhesion privacy law, and the IEEE 7012 / MyTerms
literature.

\textbf{Confidence that the residual survives: moderate (70\%).}

\subsubsection{3. CTR-LR-04 â€'' MIRAGE}\label{ctr-lr-04-uxe2-mirage}

\textbf{Claim as stated.} The observer-to-subject value gap is
market-position rent (atomisation discount + venue exclusion), not
super-additive aggregation; returns-to-scale evidence rejects raw-record
compounding; therefore redistribution requires pooling bargaining
position, not replicating aggregation scale.

\textbf{Why it MIRAGE'd.} Each plank has corpus-adjacent prior art. The
gap-as-structural-rent plank is anticipated by Bergemann-Bonatti-Gan
2022 (aggregator captures total information value via the social-data
externality) and by the Posner-Weyl / Arrieta-Ibarra monopsony account.
The flat-returns-to-scale plank is supplied by
Bajari-Chernozhukov-Hortacsu-Suzuki 2019 and the ``data is not the new
oil'' line. The pool-position plank is the standing data-union
programme. The claim rests novelty on the joint move, but the join is
closely anticipated by combining works already co-present.

\textbf{Residual novel core (survives).} The directed internal
falsifier: using flat-N returns-to-scale evidence to deny a positive
marginal product of an additional raw record, thereby showing that
data-as-labour's ``raise the wage toward marginal product via a labour
cartel'' remedy self-defeats (marginal product approximately zero), and
re-specifying the redistributive object as
venue-access/clearing-position rent rather than a wage. Stated as a
directed correction to the labour-union / propertisation programme
inside a reconstruction-bound / IEEE 7012 frame.

\textbf{Corpus tested against.} 20-item corpus plus data
returns-to-scale, data-monopsony and data-union remedy literatures,
social-data-externality surplus-capture work, and thin-market
propertisation critiques.

\textbf{Confidence that the residual survives: moderate (70\%).}

\subsubsection{4. CTR-LR-01 â€'' MIRAGE}\label{ctr-lr-01-uxe2-mirage}

\textbf{Claim as stated.} PIA measures impact to a named subject after
the fact by survey/process (ISO/PIA); PVM instead bounds disclosure in
advance and architecturally, with R a per-subject reconstruction bound
and Phi the system property gating R; the unit of analysis shifts from
subject-impact to system-property.

\textbf{Why it MIRAGE'd.} The system-property/reconstruction-bound
structure is the founding form of quantitative information flow (Smith
2009 min-entropy vulnerability is a per-secret one-guess reconstruction
bound) and differential privacy (Dwork 2006 epsilon is an ex-ante
mechanism property; Dinur-Nissim 2003 the per-record reconstruction
threshold). Separately, Wagner-Boiten 2018 already argues the PIA
tradition should be replaced by quantified system-level metrics, the
exact subject-survey-to-system-metric transition the claim names. Each
half of the move is covered; what remains is re-contextualisation.

\textbf{Residual novel core (survives, thin).} A packaging move, not a
new object: explicitly positioning a governance-facing, independently
verifiable mechanism property Phi, held two-level-distinct from the
per-subject reconstruction bound R it gates, as the single unit that a
privacy impact-assessment or valuation instrument is conducted over.
Genuinely un-verbatim only in the deliberate Phi/R
governance-versus-quantity separation offered as the assessment unit.

\textbf{Corpus tested against.} 20-item corpus plus QIF (Smith 2009), DP
(Dwork 2006), Dinur-Nissim 2003, Sweeney 2002, Wagner-Boiten 2018, and
Balle-Cherubin-Hayes 2022 as the reconstruction-bound proxy.

\textbf{Confidence that the residual survives: moderate (68\%).}

\subsubsection{5. CTR-LR-05 â€'' MIRAGE}\label{ctr-lr-05-uxe2-mirage}

\textbf{Claim as stated.} Behavioural-data-as-capital is durable only
for a finite, adversary-relative horizon set by an information-theoretic
reconstruction bound R(t) that drifts upward as adversary capability
grows; the asset has a term, not perpetuity; no prior propertisation or
valuation work indexes durability to a time-varying reconstruction
bound.

\textbf{Why it MIRAGE'd.} The headline (a durability term set by an
adversary-relative horizon that drifts upward) is substantially covered
by Mosca's cryptographic shelf-life inequality (X + Y \textgreater{} Z)
and the harvest-now-decrypt-later threat model, which already joins a
drifting-adversary horizon to a required-protection term. That is
exactly the convert-drift-into-a-maturity move the claim advertises as
unclaimed.

\textbf{Residual novel core (survives).} Substituting an
information-theoretic reconstruction bound in the Dinur-Nissim / QIF
sense for Mosca's cryptographic break-time clock, and carrying that
substituted clock into a subject-owned capital-asset valuation. Distinct
from Mosca (which prices nothing and uses a crypto break-time) and from
staleness/relevance-decay depreciation (which prices a term but on an
endogenous fixed clock). A thin combination novelty, not the structural
insight the claim advertised.

\textbf{Corpus tested against.} 20-item corpus plus Mosca inequality /
cryptographic shelf-life / HNDL, Dinur-Nissim reconstruction and
adversary-aware DP reconstruction bounds, data-value depreciation
economics, and finite-horizon data-as-capital pricing.

\textbf{Confidence that the residual survives: moderate (66\%).}

\subsection{What MIRAGE'd, and why that is a feature not a
bug}\label{what-miraged-and-why-that-is-a-feature-not-a-bug}

All five candidates were knocked down from their stated altitude. This
is the evidence that the prior-art search was adversarial rather than
confirmatory. The pattern is consistent: in every case the \emph{object}
(a reconstruction bound; an ultimatum SPE; a structural rent; a
system-property unit; a shrinking durability term) was already owned by
named prior art, and the surviving novelty is a \emph{specific
composition or re-specification} layered on top of that object.

\begin{itemize}
\tightlist
\item
  CTR-LR-02: the market-priced system property is owned by Ghosh-Roth;
  only the drift + adversary-relativity + inalienability + rent
  quadruple survives.
\item
  CTR-LR-03: the ultimatum floor and the proposer inversion are each
  owned; only the preference-free venue-exclusion parameterisation
  unifying them survives.
\item
  CTR-LR-04: rent-framing, flat returns-to-scale, and pool-for-leverage
  are each owned; only the directed self-defeat falsifier of the
  data-as-labour remedy survives.
\item
  CTR-LR-01: the reconstruction-bound object (QIF/DP) and the
  PIA-to-metric transition (Wagner-Boiten) are each owned; only the
  Phi/R governance separation as the assessment unit survives, and
  thinly.
\item
  CTR-LR-05: the shrinking adversary-relative term is owned by Mosca;
  only the IT-bound-for-crypto-clock substitution inside a capital
  valuation survives.
\end{itemize}

The disciplined WEIS contribution is therefore the \emph{conjunction} of
these surviving residues within a single instrument, not any one of the
headline claims taken alone. None of the residues should be written as
``no prior work does X''; each is a candidate tested against a named
corpus with a stated confidence, per D2.

\newpage

\section{CTR candidates: conjecture-shaped residues (assignment is the
First
Person's)}\label{ctr-candidates-conjecture-shaped-residues-assignment-is-the-first-persons}

\subsection{Conjecture candidates (residues in CTR-series
form)}\label{conjecture-candidates-residues-in-ctr-series-form}

Each surviving residual novel core is restated below as a testable
conjecture with the measurement that would settle it. These are
proposals only. \textbf{CTR-series ID assignment is the First Person's;
the IDs below are placeholders in the CTR-LR-R series and carry no
authority until the First Person ratifies them.} They are never to be
promoted to C-series here. Every proposition is a candidate against a
named corpus, not an absence claim.

\subsubsection{CTR-LR-R1 (residue of CTR-LR-02) â€''
proposed}\label{ctr-lr-r1-residue-of-ctr-lr-02-uxe2-proposed}

\textbf{Proposition.} Pricing a \emph{finite, adversary-relative,
time-drifting} reconstruction bound R(t) held as inalienable
subject-owned capital, with the observer-to-subject gap decomposed as
market-position rent, is not reducible to the
differential-privacy-markets pricing of a static budget epsilon
(Ghosh-Roth and successors).

\textbf{Measurement that would settle it.} Construct a mapping from a
Ghosh-Roth-style auction on epsilon to the PVM pricing of R(t). If the
mapping is total and value-preserving (the PVM price is recovered as a
special case by freezing drift and adversary-relativity and removing the
rent decomposition), the residue collapses and the conjecture is
refuted. If any of the four features (drift, adversary-relativity,
structural inalienability, rent-not-aggregation) has no image under the
mapping, the residue stands for that feature. Confidence the residue
stands under this test: high (86\%), speculative.

\subsubsection{CTR-LR-R2 (residue of CTR-LR-03) â€''
proposed}\label{ctr-lr-r2-residue-of-ctr-lr-03-uxe2-proposed}

\textbf{Proposition.} The measured WTA/WTP floor and the IEEE 7012
remedy are two evaluations of one preference-free SPE share function
s(proposer), where the floor equals a structural venue-exclusion
parameter h\_X (not a reservation-utility artefact), and IEEE 7012 is
the specific proposer permutation carrying h\_X to an interior share.

\textbf{Measurement that would settle it.} Specify s(proposer) with no
preference primitives and fit it to two anchor points: the observed
near-zero floor (firm-proposer arm) and any measured post-MyTerms share
(subject-proposer arm). Settled affirmatively if a single
preference-free parameter set reproduces both arms and if varying the
reservation utility alone cannot reproduce the floor (isolating venue
exclusion h\_X as the driver). Refuted if the floor is fully explained
by reservation utility without a venue-exclusion term. Confidence:
moderate (70\%), speculative.

\subsubsection{CTR-LR-R3 (residue of CTR-LR-04) â€''
proposed}\label{ctr-lr-r3-residue-of-ctr-lr-04-uxe2-proposed}

\textbf{Proposition.} Scale-replication data cooperatives fail to
redistribute value because returns to raw-record scale are flat
(marginal product of an additional record approximately zero), so the
data-as-labour ``wage toward marginal product via a cartel'' remedy
self-defeats, and only pooling venue/clearing-position rent
redistributes.

\textbf{Measurement that would settle it.} Estimate the elasticity of
downstream model or auction value with respect to N (records) holding
the clearing venue fixed, versus its elasticity with respect to
bargaining/venue access holding N fixed. Settled affirmatively if the
N-elasticity is statistically indistinguishable from zero over the
relevant range while the venue-access elasticity is strictly positive.
Refuted if raw-record N carries a positive, exploitable marginal product
that a cooperative could monetise as a wage. Confidence: moderate
(70\%), speculative.

\subsubsection{CTR-LR-R4 (residue of CTR-LR-01) â€''
proposed}\label{ctr-lr-r4-residue-of-ctr-lr-01-uxe2-proposed}

\textbf{Proposition.} A governance-facing, independently verifiable
mechanism property Phi, held two-level-distinct from the per-subject
reconstruction bound R it gates, can serve as the single unit over which
a privacy impact-assessment or valuation instrument is conducted, and
this two-level separation is not already carried by QIF/DP (which supply
the object but not the instrument framing) nor by Wagner-Boiten (which
supplies the instrument framing but not the Phi-gates-R unit).

\textbf{Measurement that would settle it.} Attempt to instantiate the
PIA/valuation instrument using QIF/DP objects alone (no explicit Phi/R
governance split) and, separately, using Wagner-Boiten metrics alone.
Settled affirmatively if neither instantiation reproduces the
auditability property that the explicit Phi/R separation provides (for
example, a third-party attestation of Phi without disclosure of R).
Refuted if either prior instrument already yields that attestation. This
residue is thin; confidence: moderate (68\%), speculative.

\subsubsection{CTR-LR-R5 (residue of CTR-LR-05) â€''
proposed}\label{ctr-lr-r5-residue-of-ctr-lr-05-uxe2-proposed}

\textbf{Proposition.} Substituting an information-theoretic
reconstruction bound R(t) (Dinur-Nissim / QIF sense) for Mosca's
cryptographic break-time clock, and carrying it into a subject-owned
capital-asset valuation, yields a maturity/term that neither Mosca
(prices nothing, crypto clock) nor staleness-decay depreciation (prices
a term, fixed endogenous clock) produces.

\textbf{Measurement that would settle it.} Compute the durability term
three ways for the same behavioural-data asset: (a) Mosca cryptographic
shelf-life, (b) staleness-decay depreciation, (c) IT reconstruction
bound R(t). Settled affirmatively if the R(t) term diverges materially
from both (a) and (b) under a growing-adversary schedule (for example
the term shortens with adversary compute in a way the fixed-clock model
cannot express and the crypto-clock model expresses only for
confidentiality, not valuation). Refuted if all three coincide up to
relabelling. Confidence: moderate (66\%), speculative.

\textbf{Note on scope.} None of the above is stated as ``no prior work
does X.'' Each is a candidate proposition bound to the corpus and web
scope named in its parent verdict, with a confidence label, per
discipline D2. Promotion, renumbering, or rejection is reserved to the
First Person.

\newpage

\section{CTR-LR-02b: DP-refinements stress test of the R(t) residue (LM2
resolution)}\label{ctr-lr-02b-dp-refinements-stress-test-of-the-rt-residue-lm2-resolution}

\subsection{Verdict:
distinction\_holds\_but\_phrasing\_wrong}\label{verdict-distinction_holds_but_phrasing_wrong}

\textbf{Residual after stress:} The R(t) residue survives the LM2 stress
test, but sharpened and narrowed rather than left at its LR-02 form.
Honest confidence: \textasciitilde85\% (essentially steady on the
empty-cell fact, downgraded on its reach). The high-confidence,
scope-bounded part is that within the surveyed DP-refinements strand no
work indexes a fixed archive's per-subject reconstruction residual to
adversary-capability growth along calendar time (CTR-LR-02b-ii/iii,
high). What the stress test removes is the overreach: the claim can no
longer be stated as `categorically distinct' or grounded on `static,
adversary-independent' DP. The surviving form is (i) a CLOCK distinction
(query/release-composition vs calendar-time adversary capability), (ii)
restricted to the INFORMATIONAL-capability reading (accumulating
linkage/side-prior corpus, not compute, which orthodox DP saturates),
and (iii) explicitly current-formalization-contingent, since
adversary-aware DP already carries capability as a static parameter and
computational DP already gestures at the same erosion. Publishable as
sharpened novelty in that narrowed informational form; not as an
in-principle impossibility result.

\subsection{What differential privacy DOES do (do not
understate)}\label{what-differential-privacy-does-do-do-not-understate}

Adaptive, Rényi, concentrated, Gaussian and composition DP do model an
adversary: the guarantee is calibrated against a worst-case adversary of
unbounded computation and arbitrary auxiliary knowledge, entering as a
supremum (Mironov, 2017; Bun and Steinke, 2016; Dong, Roth and Su,
2022). They also carry a genuine temporal axis: privacy loss epsilon is
accounted as it accumulates across composed mechanisms on a
query/release-composition clock (Kairouz, Oh and Viswanath, 2015; Dwork,
Rothblum and Vadhan, 2010), and continual-observation DP tracks a stream
of T published outputs (Dwork, Naor, Pitassi and Rothblum, 2010).
Adversary-aware refinements go further and parametrise which fixed
adversary class the guarantee is calibrated against (Cummings et al.,
2024; Swanberg et al., 2025).

\subsection{What differential privacy does NOT do (relative to
R(t))}\label{what-differential-privacy-does-not-do-relative-to-rt}

None of these index a fixed, never-re-queried archive's per-subject
reconstruction residual to the adversary's informational capability
(accumulating linkage corpus, side priors, informedness) growing along a
calendar clock while nothing is re-released. Their temporal axis is
added queries or releases (k, or stream-length T) against a held-fixed
worst-case adversary, which is the inverse polarity to a frozen release
aging, and where DP does run a wall clock (age-dependent and
temporally-discounted DP; Valavi et al., 2022) the guarantee strengthens
as data ages through relevance decay. Adversary capability therefore
enters only as a static supremum or a fixed class parameter, never as a
variable integrated along calendar time; the one caveat is that this
exactness is proven for information-theoretic DP, with computational DP
a separate, out-of-scope literature.

\subsection{Sub-claim adjudication}\label{sub-claim-adjudication}

\begin{itemize}
\item
  \textbf{CTR-LR-02b-i: PARTIAL} (moderate-high) - DP composition and
  the information-theoretic DP family (RDP/zCDP/GDP/f-DP,
  advanced/optimal/continual-observation composition) provably run a
  query/release clock â€'' free variable k, adversary held at worst-case
  sup\_C â€'' and therefore cannot express a fixed release's
  reconstruction residual rising with adversary capability over calendar
  time; R(t) is a distinct clock (free variable calendar-time t through
  a growing adversary C(t), single frozen release k=1), and no DP work
  in the searched corpus formalizes it. This distinctness is
  contingent-on-current-formalization, not categorical-in-principle:
  computational DP already recognizes (informally) that a fixed
  release's guarantee erodes as adversary compute grows over calendar
  time â€'' the same mechanism, unformalized â€'' and adversary-aware DP
  bounds (ATTAXONOMY / 2507.08158) already carry adversary capability C
  as an explicit static parameter, so a calendar-indexed C(t) is a
  natural, non-foreclosed bridge rather than a mathematical
  impossibility.
\item
  \textbf{CTR-LR-02b-ii: HOLDS} (high) - Within the enumerated
  DP-refinements strand â€'' all information-theoretic (RDP, zCDP/CDP,
  advanced/optimal/concurrent composition, continual-observation DP,
  f-DP/GDP, adversary-aware accounting) plus the time-aware DP families
  (temporal-correlation TPL/ConTPL/CSDP, age-dependent,
  temporally-discounted, everlasting) â€'' no work indexes a fixed
  release's residual reconstruction/inference risk to
  adversary-CAPABILITY growth over calendar time. Each runs one of three
  non-R(t) clocks: (1) composition/query accounting with a fixed
  worst-case adversary, (2) release-stream length T of inverted polarity
  with a fixed adversary, or (3) a static adversary
  prior/baseline-success parameter with no time variable. The
  temporally-structured members (incl.~Cao's TPL) run the stream clock
  against a fixed adversary; the wall-clock members (age-dependent,
  temporally-discounted) run relevance/age decay of the OPPOSITE
  polarity (guarantee strengthens as data ages). The structural reason
  is exact for this strand: every member is information-theoretic and
  robust to arbitrary auxiliary knowledge and unbounded compute, so
  adversary capability enters only as a supremum, never a calendar
  index. Out-of-scope caveat: computational DP and the crypto
  harvest-now-decrypt-later framing could in principle carry a
  capability-growth clock, but they are a different literature already
  excluded and do not formally index epsilon to calendar time.
\item
  \textbf{CTR-LR-02b-iii: PARTIAL} (high) - Over the
  DP-refinements/DP-markets strand (\textasciitilde40 works surveyed),
  the joint cell C = A x B x C is empty and no member co-occupies even
  two of its three axes, where A = an information-theoretic per-subject
  reconstruction bound on a FIXED (non-growing) archive; B = a
  Mosca-STRUCTURED calendar-time clock along which the adversary's
  AUXILIARY/INFORMATIONAL capability (linkage corpus, side priors,
  informedness â€'' NOT computational power, which axis A renders inert
  and which unbounded-adversary DP already saturates) grows, so the
  frozen archive's per-subject reconstruction residual rises
  monotonically with t though nothing is re-released â€'' explicitly
  neither a crypto break-time nor a DP composition/continual-release
  budget; and C = the residual held and priced as inalienable
  subject-owned CAPITAL (a stock) rather than a per-query privacy-loss
  FLOW. The clock distinction is real and exact against every surveyed
  clock; the contribution survives as novelty only in this narrower,
  informational-capability form.
\end{itemize}

\subsection{Applied to the draft}\label{applied-to-the-draft}

\begin{itemize}
\item
  S1 fix: The differential-privacy budget is not adversary-independent;
  it is calibrated against a worst-case adversary of unbounded
  computation and arbitrary auxiliary knowledge, which enters the
  guarantee as a supremum (Mironov, 2017; Bun and Steinke, 2016; Dong,
  Roth and Su, 2022). What that literature holds fixed is not the
  adversary but the clock: privacy loss is accounted along a
  query/release-composition axis, with epsilon degrading as mechanisms
  are composed (Kairouz, Oh and Viswanath, 2015; Dwork, Rothblum and
  Vadhan, 2010) or as a stream of T outputs is published under continual
  observation (Dwork, Naor, Pitassi and Rothblum, 2010), the adversary's
  capability being held at that worst-case supremum throughout.
  {[}Placement: replaces the S1 sentence that currently characterises
  the DP-markets budget as `static, adversary-independent', i.e.~the
  positioning sentence that opens the DP-refinements comparison.{]}
\item
  S2 fix: R(t) runs a different clock. It indexes a single fixed
  release, never re-queried, whose per-subject reconstruction residual
  rises as the adversary's informational capability, namely its
  accumulating linkage corpus and side priors rather than its
  computational power (which differential privacy already saturates),
  grows along calendar time. No work in the surveyed DP-refinements
  strand runs this clock: composition and its Rényi, concentrated,
  Gaussian and continual-observation refinements all move along added
  queries or releases against a fixed worst-case adversary, and the
  wall-clock DP families, being age-dependent and temporally-discounted
  DP together with the data-value decay of Valavi et al.~(2022), run
  relevance decay of the opposite polarity in which the guarantee
  strengthens as data ages. We state this as a clock distinction and not
  a categorical barrier: adversary-aware DP already carries adversary
  capability as an explicit static parameter (Cummings et al., 2024;
  Swanberg et al., 2025) and computational DP already informally
  recognises that a fixed release erodes as adversary compute grows, so
  a calendar-indexed capability clock is a natural but
  as-yet-unformalised extension rather than a mathematical
  impossibility. {[}Placement: replaces the S2 sentence that restates
  `static, adversary-independent' as the ground for the R(t)
  distinction, i.e.~the contrast sentence that asserts R(t)'s novelty
  against the strand.{]}
\end{itemize}

\subsection{Ledger paragraph}\label{ledger-paragraph}

LM2 resolved as distinction\_holds\_but\_phrasing\_wrong. The DP referee
is right that
adaptive/Rényi/concentrated/Gaussian/composition/continual-observation
DP model adversaries (a worst-case, unbounded-compute,
arbitrary-auxiliary-knowledge adversary entering as a supremum) and
carry a temporal axis (privacy loss accounted along query/release
composition, or a stream of T releases under continual observation), and
adversary-aware DP even parametrises the fixed adversary class; so
`static, adversary-independent' is a genuine mischaracterisation and
must be replaced. But the R(t) distinction does not collapse
(adjudicated CTR-LR-02b-i PARTIAL, -ii HOLDS, -iii PARTIAL): every
strand member runs one of three non-R(t) clocks (composition/query
accounting against a fixed worst-case adversary; release-stream length T
of inverted polarity; or a static adversary-class parameter with no time
variable), and the wall-clock DP families (age-dependent,
temporally-discounted, Valavi et al.~2022) run relevance decay of the
OPPOSITE polarity in which the guarantee strengthens as data ages. The
residue is re-stated as a CLOCK distinction (query/release-composition
vs calendar-time adversary capability) and narrowed to the
informational-capability reading forced by CTR-LR-02b-iii (accumulating
linkage/side priors, not compute, which DP already saturates). Two
honesty constraints are carried into the prose: the distinction is
current-formalization-contingent, not categorical (computational DP and
adversary-aware bounds are non-foreclosed bridges), and every negative
claim is scoped to the surveyed DP-refinements strand per D2. S1/S2
drop-in replacements written engaging Mironov 2017, Bun-Steinke 2016,
Dong-Roth-Su 2022, Kairouz-Oh-Viswanath 2015, Dwork-Rothblum-Vadhan
2010, Dwork-Naor-Pitassi-Rothblum 2010, Cummings et al.~2024, Swanberg
et al.~2025 and Valavi et al.~2022 by author-year. Residual after
stress: \textasciitilde85\%, sharper but narrower than the 86\% LR-02
figure.

\newpage

\section{Bibliography: challenge corpus with
provenance}\label{bibliography-challenge-corpus-with-provenance}

\subsection{Bibliography (WEIS-submission
ready)}\label{bibliography-weis-submission-ready}

Grouped by strand. Every item below carries
\texttt{provenance\_confidence\ =\ verified\_from\_search}. \textbf{No
item is flagged \texttt{uncertain}; no first-person provenance
verification is outstanding for this corpus.} Standards items list both
the cited edition and its known revision for the First Person to select
the citation edition at submission.

\subsubsection{Strand A â€'' Propertisation and property rights in
personal data
(legal)}\label{strand-a-uxe2-propertisation-and-property-rights-in-personal-data-legal}

\begin{enumerate}
\def\labelenumi{\arabic{enumi}.}
\tightlist
\item
  Kenneth C. Laudon, ``Markets and Privacy,'' \emph{Communications of
  the ACM}, Vol. 39, No.~9 (September 1996), pp.~92-104.
  https://doi.org/10.1145/234215.234476 (verified\_from\_search)
\item
  Pamela Samuelson, ``Privacy As Intellectual Property?'', 52
  \emph{Stanford Law Review} 1125-1173 (2000).
  https://www.jstor.org/stable/1229511 (verified\_from\_search)
\item
  Paul M. Schwartz, ``Property, Privacy, and Personal Data,'' 117
  \emph{Harvard Law Review} 2056 (2004).
  https://papers.ssrn.com/sol3/papers.cfm?abstract\_id=721642
  (verified\_from\_search)
\item
  Nadezhda Purtova, ``The illusion of personal data as no one's
  property,'' \emph{Law, Innovation and Technology} 7(1) (2015): 83-111;
  see also N. Purtova, \emph{Property Rights in Personal Data: A
  European Perspective} (Kluwer Law International, 2011/2012) and ``Do
  Property Rights in Personal Data Make Sense after the Big Data Turn?
  Individual Control and Transparency'' (2017).
  https://doi.org/10.1080/17579961.2015.1052646 (verified\_from\_search)
\item
  J.E.J. (Corien) Prins, ``Property and Privacy: European Perspectives
  and the Commodification of our Identity,'' in L. Guibault and P.B.
  Hugenholtz (eds), \emph{The Future of the Public Domain: Identifying
  the Commons in Information Law} (Information Law Series No.~16,
  pp.~223-257), Kluwer Law International, 2006.
  https://papers.ssrn.com/sol3/papers.cfm?abstract\_id=929668
  (verified\_from\_search)
\end{enumerate}

\subsubsection{Strand B â€'' Data-as-labour and market-structure
remedies}\label{strand-b-uxe2-data-as-labour-and-market-structure-remedies}

\begin{enumerate}
\def\labelenumi{\arabic{enumi}.}
\setcounter{enumi}{5}
\tightlist
\item
  Imanol Arrieta-Ibarra, Leonard Goff, Diego Jimenez-Hernandez, Jaron
  Lanier, and E. Glen Weyl, ``Should We Treat Data as Labor? Moving
  beyond `Free','' \emph{AEA Papers and Proceedings} 108 (2018): 38-42.
  https://doi.org/10.1257/pandp.20181003 (verified\_from\_search)
\item
  Eric A. Posner and E. Glen Weyl, ``Data as Labor: Valuing
  Contributions to the Digital Economy,'' Chapter 5 in \emph{Radical
  Markets: Uprooting Capitalism and Democracy for a Just Society},
  pp.~205-249, Princeton University Press, 2018.
  https://doi.org/10.1515/9780691196978-009 (verified\_from\_search)
\item
  Sylvie Delacroix and Neil D. Lawrence, ``Bottom-up data Trusts:
  disturbing the `one size fits all' approach to data governance,''
  \emph{International Data Privacy Law} 9(4) (2019): 236-252.
  https://doi.org/10.1093/idpl/ipz014 (verified\_from\_search)
\end{enumerate}

\subsubsection{Strand C â€'' Empirical valuation of privacy (WTA/WTP,
market
surveys)}\label{strand-c-uxe2-empirical-valuation-of-privacy-wtawtp-market-surveys}

\begin{enumerate}
\def\labelenumi{\arabic{enumi}.}
\setcounter{enumi}{8}
\tightlist
\item
  Alessandro Acquisti, Curtis Taylor, and Liad Wagman, ``The Economics
  of Privacy,'' \emph{Journal of Economic Literature} 54(2) (2016):
  442-492. https://doi.org/10.1257/jel.54.2.442 (verified\_from\_search)
\item
  Alessandro Acquisti, Leslie K. John, and George Loewenstein, ``What Is
  Privacy Worth?'', \emph{The Journal of Legal Studies} 42(2) (2013):
  249-274. https://doi.org/10.1086/671754 (verified\_from\_search)
\item
  Avinash Collis, Alex Moehring, Ananya Sen, and Alessandro Acquisti,
  ``Information Frictions and Heterogeneity in Valuations of Personal
  Data,'' 21st Workshop on the Economics of Information Security (WEIS
  2022). Median WTA to share entirety of Facebook data = \$750 (YouGov
  nationally representative sample); \$1,000 (Data Dividend Project
  sample); bimodal distribution clustered below \$250 and above
  \$10,000.
  https://weis2022.econinfosec.org/wp-content/uploads/sites/10/2022/06/weis22-collis.pdf
  (verified\_from\_search)
\item
  Alastair R. Beresford, Dorothea Kuebler, and Soeren Preibusch,
  ``Unwillingness to pay for privacy: A field experiment,''
  \emph{Economics Letters} 117(1) (2012): 25-27.
  https://doi.org/10.1016/j.econlet.2012.04.077 (verified\_from\_search)
\item
  Jens Grossklags and Alessandro Acquisti, ``When 25 Cents is Too Much:
  An Experiment on Willingness-To-Sell and Willingness-To-Protect
  Personal Information,'' Sixth Workshop on the Economics of Information
  Security (WEIS 2007), Carnegie Mellon University, June 2007.
  https://econinfosec.org/archive/weis2007/papers/66.pdf
  (verified\_from\_search)
\item
  Sarah Spiekermann, Alessandro Acquisti, Rainer Boehme, and Kai-Lung
  Hui, ``The challenges of personal data markets and privacy,''
  \emph{Electronic Markets} 25(2) (2015): 161-167.
  https://doi.org/10.1007/s12525-015-0191-0 (verified\_from\_search)
\end{enumerate}

\subsubsection{Strand D â€'' Security economics (incentive
framing)}\label{strand-d-uxe2-security-economics-incentive-framing}

\begin{enumerate}
\def\labelenumi{\arabic{enumi}.}
\setcounter{enumi}{14}
\tightlist
\item
  Ross Anderson, ``Why Information Security is Hard - An Economic
  Perspective,'' \emph{Proceedings of the 17th Annual Computer Security
  Applications Conference (ACSAC 2001)}, New Orleans, LA, 10-14 December
  2001, pp.~358-365, IEEE Computer Society.
  https://doi.org/10.1109/ACSAC.2001.991552 (draws on G. Akerlof, ``The
  Market for Lemons,'' \emph{QJE} 84(3) (1970): 488-500).
  (verified\_from\_search)
\end{enumerate}

\subsubsection{Strand E â€'' Formal privacy metrics and system
properties}\label{strand-e-uxe2-formal-privacy-metrics-and-system-properties}

\begin{enumerate}
\def\labelenumi{\arabic{enumi}.}
\setcounter{enumi}{15}
\tightlist
\item
  Cynthia Dwork, ``Differential Privacy,'' in \emph{Automata, Languages
  and Programming (ICALP 2006)}, Part II, LNCS vol.~4052, pp.~1-12,
  Springer, 2006. https://doi.org/10.1007/11787006\_1
  (verified\_from\_search)
\item
  Latanya Sweeney, ``k-Anonymity: A Model for Protecting Privacy,''
  \emph{International Journal of Uncertainty, Fuzziness and
  Knowledge-Based Systems} 10(5) (2002): 557-570.
  https://doi.org/10.1142/S0218488502001648 (verified\_from\_search)
\item
  Geoffrey Smith, ``On the Foundations of Quantitative Information
  Flow,'' in \emph{Foundations of Software Science and Computational
  Structures (FOSSACS 2009)}, LNCS vol.~5504, pp.~288-302, Springer,
  2009. https://doi.org/10.1007/978-3-642-00596-1\_21
  (verified\_from\_search)
\end{enumerate}

\subsubsection{Strand F â€'' Standards (process and
framework)}\label{strand-f-uxe2-standards-process-and-framework}

\begin{enumerate}
\def\labelenumi{\arabic{enumi}.}
\setcounter{enumi}{18}
\tightlist
\item
  ISO/IEC 29134:2017, \emph{Information technology - Security techniques
  - Guidelines for privacy impact assessment}, 1st ed., ISO/IEC, Geneva,
  2017 (revised by ISO/IEC 29134:2023).
  https://www.iso.org/standard/62289.html (verified\_from\_search).
  First-person note: select 2017 or 2023 edition at submission.
\item
  ISO/IEC 29100:2011, \emph{Information technology - Security techniques
  - Privacy framework}, ISO/IEC (second edition ISO/IEC 29100:2024).
  https://www.iso.org/standard/45123.html (verified\_from\_search).
  First-person note: select 2011 or 2024 edition at submission.
\end{enumerate}

\subsubsection{External works cited in adjudication (not in the 20-item
corpus, listed for the reference
apparatus)}\label{external-works-cited-in-adjudication-not-in-the-20-item-corpus-listed-for-the-reference-apparatus}

These were surfaced by the adjudicators' web scope and are load-bearing
for the MIRAGE verdicts; the First Person may wish to cite them in the
related-work section. Provenance: adjudicator web scope, July 2026, not
independently re-verified here (flag for first-person verification
before citation).

\begin{itemize}
\tightlist
\item
  Arpita Ghosh and Aaron Roth, ``Selling Privacy at Auction,'' EC 2011
  (DP-markets strand; occupant of the \{system\_property x
  market\_pricing\} cell). Verification pending.
\item
  Dinesh Bergemann, Alessandro Bonatti, and Tan Gan, ``The Economics of
  Social Data,'' RAND Journal of Economics, 2022 (aggregator surplus
  capture, social-data externality). Verification pending.
\item
  Patrick Bajari, Victor Chernozhukov, Ali Hortacsu, and Junichi Suzuki,
  ``The Impact of Big Data on Firm Performance,'' AEA P\&P, 2019 (flat
  returns-to-scale evidence). Verification pending.
\item
  Ben Wagner and Eerke Boiten, ``Privacy Risk Assessment: From Art to
  Science, By Metrics,'' DPM 2018 (PIA-to-metric transition).
  Verification pending.
\item
  Irit Dinur and Kobbi Nissim, ``Revealing Information while Preserving
  Privacy,'' PODS 2003 (per-record reconstruction threshold).
  Verification pending.
\item
  Michele Mosca, cryptographic shelf-life inequality (X + Y
  \textgreater{} Z) and harvest-now-decrypt-later threat model (IEEE
  S\&P 2018; FEDS 2025-093). Verification pending.
\item
  IEEE 7012 / Customer Terms (MyTerms), ProjectVRM / Doc Searls,
  myterms.info (proposer-inversion standard). Verification pending.
\end{itemize}

\end{document}
